\section{ResNet}

La ResNet (Residual Network) introduce le \textit{residual connections} per facilitare
l'addestramento di reti molto profonde. L'idea è che gli strati convoluzionali non
debbano necessariamente apprendere direttamente la funzione desiderata $H(x)$, ma
possano apprendere una funzione residua $F(x)$ rispetto all'input:

\[
F(x) = H(x) - x
\]

Il risultato del blocco viene quindi ottenuto come:

\[
H(x) = F(x) + x
\]

L'input $x$ viene quindi mantenuto attraverso una \textit{shortcut connection} e
sommato all'output della trasformazione residua. Questa struttura facilita
l'ottimizzazione delle reti profonde, poiché, nel caso in cui la trasformazione
ottimale fosse vicina all'identità, il blocco deve imparare una correzione rispetto
all'input invece dell'intera trasformazione. Il paper originale introduce questa
formulazione come principio fondamentale della residual learning. 

\subsection{Bottleneck}

ResNet-50 utilizza un particolare tipo di residual block chiamato
\textit{bottleneck block}. Il blocco è costituito da tre convoluzioni:

\[
1\times1 \rightarrow 3\times3 \rightarrow 1\times1
\]

La prima convoluzione $1\times1$ riduce il numero di canali della rappresentazione.
La convoluzione $3\times3$, che è la principale convoluzione spaziale del blocco,
opera quindi su un numero ridotto di canali, diminuendo il costo computazionale.
Infine, la seconda convoluzione $1\times1$ aumenta nuovamente il numero di canali,
riportandolo alla dimensione necessaria per effettuare la somma con la shortcut.

Ad esempio, un blocco può trasformare una rappresentazione:

\[
256 \rightarrow 64 \rightarrow 64 \rightarrow 256
\]

dove il passaggio a 64 canali costituisce il \textit{bottleneck}.

La struttura completa del blocco può essere rappresentata come:

\[
x
\rightarrow
1\times1
\rightarrow
BN
\rightarrow
ReLU
\rightarrow
3\times3
\rightarrow
BN
\rightarrow
ReLU
\rightarrow
1\times1
\rightarrow
BN
\]

e, in parallelo, l'input viene trasmesso attraverso la shortcut. I due rami vengono
poi sommati e viene applicata una ReLU finale:

\[
y = ReLU(F(x) + x)
\]

Quando le dimensioni dell'input e dell'output non coincidono, la shortcut viene
trasformata tramite una convoluzione $1\times1$ per adattarne numero di canali e,
se necessario, risoluzione spaziale.

Il bottleneck è stato introdotto nel paper originale proprio per rendere più
efficiente la costruzione di reti molto profonde: le convoluzioni $1\times1$
riducano e poi aumentano le dimensioni, lasciando alla convoluzione $3\times3$ una
rappresentazione più piccola su cui operare. 

\subsubsection{Batch Normalization}

La \textit{Batch Normalization} (BN) normalizza le attivazioni prodotte da una
convoluzione utilizzando la media e la varianza calcolate sul mini-batch. 
Successivamente, le attivazioni normalizzate vengono trasformate tramite due
parametri apprendibili, che permettono alla rete di adattarne nuovamente scala e
traslazione.

In ResNet, la Batch Normalization viene applicata dopo ogni convoluzione e prima
della funzione di attivazione ReLU. La sua funzione è rendere più stabile
l'addestramento della rete e facilitare l'ottimizzazione dei numerosi strati
convoluzionali.


\subsection{Struttura di ResNet-50}

ResNet-50 è costruita come una sequenza di quattro stage principali, ciascuno
composto da più Bottleneck block. La configurazione di ResNet-50 utilizza
rispettivamente:

\[
[3,\,4,\,6,\,3]
\]

Bottleneck nei quattro stage \texttt{conv2\_x}, \texttt{conv3\_x},
\texttt{conv4\_x} e \texttt{conv5\_x}.

La struttura può essere riassunta come:

\[
\text{Input}
\rightarrow
7\times7\ \text{Conv}
\rightarrow
\text{MaxPool}
\rightarrow
\text{conv2\_x}
\rightarrow
\text{conv3\_x}
\rightarrow
\text{conv4\_x}
\rightarrow
\text{conv5\_x}
\]

Per un input RGB di dimensione $224\times224$, le principali dimensioni delle
feature map sono:

\[
224\times224
\rightarrow
112\times112
\rightarrow
56\times56
\rightarrow
28\times28
\rightarrow
14\times14
\rightarrow
7\times7
\]

mentre il numero di canali aumenta progressivamente:

\[
64
\rightarrow
256
\rightarrow
512
\rightarrow
1024
\rightarrow
2048.
\]

Il primo layer utilizza una convoluzione $7\times7$ con stride $2$, seguita da
un max pooling con stride $2$. I successivi cambiamenti di risoluzione avvengono
all'inizio degli stage \texttt{conv3\_x}, \texttt{conv4\_x} e \texttt{conv5\_x},
dove il primo Bottleneck utilizza stride $2$. Nel codice, quindi, il primo
Bottleneck di ogni nuovo stage è anche quello che adatta la dimensione della
shortcut quando necessario.



\subsubsection{Global Average Pooling}

Al termine della rete, la ResNet originale utilizza un \textit{Global Average
Pooling} (GAP) per trasformare la feature map finale in un vettore compatto.

Nel caso di ResNet-50, la feature map finale ha dimensione:

\[
7\times7\times2048.
\]

Il Global Average Pooling calcola, per ogni canale, la media dei valori presenti
su tutte le posizioni spaziali:

\[
y_c = \frac{1}{H W}
\sum_{i=1}^{H}\sum_{j=1}^{W} x_{c,i,j}.
\]

La dimensione spaziale viene quindi ridotta da $7\times7$ a $1\times1$,
ottenendo:

\[
7\times7\times2048
\rightarrow
1\times1\times2048.
\]

Successivamente il tensore viene appiattito e passato al layer fully connected
che produce le predizioni finali.

Il vantaggio del Global Average Pooling è che permette di riassumere l'intera
mappa spaziale di ciascun canale in un singolo valore, evitando l'utilizzo di
una grande quantità di parametri tipica di un fully connected applicato
direttamente alla feature map.
Il GAP non verrà usato nel caso in cui ResNet debba funzioanr ecome backbone.